\documentclass[../../main.tex]{subfiles}

\begin{document}

\section{Nagata-Smirnov 度量化定理}\label{section:ch06-nagata-smirnov}

现在我们来证明空间 $X$ 是正则的且存在可数局部有限基就等价于 $X$ 可度量化.

这些条件蕴涵可度量性的证明, 非常接近于我们给出的 Urysohn 度量化定理的第二个证明. 我们在那个证明中构造了一个从 $X$ 到 $\mathbb{R}^\omega$ 的映射, 它是关于 $\mathbb{R}^\omega$ 的一致度量 $\bar{\rho}$ 的一个嵌入. 这就使我们想起这个证明的主要步骤. 证明的第一步是证明每一个具有可数基的正则空间是正规的. 第二步是构造 $X$ 上的一个可数的实值函数族 $\{f_n\}$, 使它能分离点和闭子集. 第三步是利用函数 $f_n$ 定义一个映射, 把 $X$ 嵌入到积空间 $\mathbb{R}^\omega$ 中. 第四步是说明, 若对任意 $x\in X$, 都有 $f_n(x)\leqslant 1/n$, 那么这个映射实际上是从 $X$ 到度量空间 $(\mathbb{R}^\omega,\bar{\rho})$ 的一个嵌入.

为了证明一般的度量化定理, 这些步骤的每一步都要加以推广. 第一, 我们证明具有可数局部有限基的正则空间是正规的. 第二, 在 $X$ 上构造某个实值函数族 $\{f_n\}$. 第三, 用这些函数把 $X$ 嵌入到积空间 $\mathbb{R}^J$ 中(对某一个 $J$). 第四步是说明, 若函数 $f_\alpha$ 充分小, 那么实际上它将 $X$ 嵌入到度量空间 $(\mathbb{R}^J,\bar{\rho})$ 中.

在证明之前, 我们回顾一下在前面习题中引入的 $G_\delta$ 集的概念.

\begin{definition}\label{definition:ch06-g-delta}
如果空间 $X$ 的一个子集 $A$ 等于 $X$ 的可数个开子集的交, 则称 $A$ 为 $X$ 中的一个 $G_\delta$ 集($G_\delta$ set).
\end{definition}

\begin{examples}
显然, $X$ 的每一个开子集是一个 $G_\delta$ 集. 在第一可数的 Hausdorff 空间中, 每一个单点集是一个 $G_\delta$ 集. 可以验证, $\bar{S}_\Omega$ 的单点子集 $\{\Omega\}$ 不是 $G_\delta$ 集.
\end{examples}

\begin{examples}
在度量空间 $X$ 中, 每一个闭集是一个 $G_\delta$ 集. 给定 $A\subset X$, 用 $U(A,\varepsilon)$ 表示 $A$ 的 $\varepsilon$-邻域. 可以验证, 若 $A$ 是闭集, 则有
\begin{align*}
A=\bigcap_{n\in\mathbb{Z}_+}U(A,1/n).
\end{align*}
\end{examples}

\begin{lemma}\label{lemma:ch06-regular-countable-locally-finite-normal}
设 $X$ 是一个正则空间并且具有一个可数局部有限基 $\mathcal{B}$, 则 $X$ 是正规的, 并且 $X$ 中每一个闭集都是 $X$ 的一个 $G_\delta$ 集.
\end{lemma}

\begin{proof}
第一步. 设 $W$ 是 $X$ 的一个开集, 我们将证明存在 $X$ 的开集的一个可数族 $\{U_n\}$, 使得
\begin{align*}
W=\bigcup U_n=\bigcup\overline{U_n}.
\end{align*}

因为 $\mathcal{B}$ 是 $X$ 的一个可数局部有限基, 所以可以写成 $\mathcal{B}=\bigcup\mathcal{B}_n$, 其中每一个族 $\mathcal{B}_n$ 都是局部有限的. 设 $\mathcal{C}_n$ 是 $\mathcal{B}$ 中满足 $B\in\mathcal{B}_n$ 并且 $\overline{B}\subset W$ 的那些元素 $B$ 构成的族. 那么 $\mathcal{C}_n$ 是局部有限的, 因为它是 $\mathcal{B}_n$ 的子集的一个族. 定义
\begin{align*}
U_n=\bigcup_{B\in\mathcal{C}_n}B,
\end{align*}
那么 $U_n$ 是开集, 并且根据\cref{lemma:ch06-locally-finite-properties},
\begin{align*}
\overline{U_n}=\bigcup_{B\in\mathcal{C}_n}\overline{B}.
\end{align*}
因此 $\overline{U_n}\subset W$, 从而
\begin{align*}
\bigcup U_n\subset\bigcup\overline{U_n}\subset W.
\end{align*}
我们断言等号成立. 给定 $x\in W$, 由正则性可见, 存在 $B\in\mathcal{B}$ 使得 $x\in B$ 和 $\overline{B}\subset W$. 现在, 对某一个 $n$ 有 $B\in\mathcal{B}_n$. 根据定义可见 $B\in\mathcal{C}_n$, 所以 $x\in U_n$. 于是 $W\subset\bigcup U_n$.

第二步. 我们证明 $X$ 的每一个闭集 $C$ 都是 $G_\delta$ 集. 给定闭集 $C$, 令 $W=X-C$. 由第一步知, 存在 $X$ 中的集合 $U_n$, 使得 $W=\bigcup\overline{U_n}$. 于是
\begin{align*}
C=\bigcap(X-\overline{U_n}),
\end{align*}
即 $C$ 等于 $X$ 中可数个开集的交.

第三步. 证明 $X$ 是正规的. 设 $C$ 和 $D$ 为 $X$ 中两个无交的闭集. 对于开集 $X-D$ 应用第一步构造一个可数开集族 $\{U_n\}$, 使得 $\bigcup U_n=\bigcup\overline{U_n}=X-D$. 于是 $\{U_n\}$ 覆盖 $C$, 并且每一个 $\overline{U_n}$ 都与 $D$ 无交. 类似地, 存在 $D$ 的可数开覆盖 $\{V_n\}$, 并且每一个 $V_n$ 的闭包也与 $C$ 无交.

现在我们再来回顾关于具有可数基的正则空间是正规的空间的证明(定理 32.1), 这里可以一字不变地重新陈述那个证明. 令
\begin{align*}
U_n'=U_n-\bigcup_{i=1}^n\overline{V_i},\quad V_n'=V_n-\bigcup_{i=1}^n\overline{U_i}.
\end{align*}
则集合
\begin{align*}
U'=\bigcup_{n\in\mathbb{Z}_+}U_n',\quad V'=\bigcup_{n\in\mathbb{Z}_+}V_n'
\end{align*}
是分别包含 $C$ 和 $D$ 的无交的开集.
\end{proof}

\begin{lemma}\label{lemma:ch06-closed-g-delta-function}
设 $X$ 是一个正规空间, $A$ 是 $X$ 的一个闭的 $G_\delta$ 子集. 则存在连续函数 $f\colon X\to[0,1]$, 使得当 $x\in A$ 时有 $f(x)=0$; 当 $x\notin A$ 时有 $f(x)>0$.
\end{lemma}

\begin{proof}
在第 33 节中曾将此命题作为练习给出, 在这里我们给出证明. 设 $A$ 为可数个开集 $U_n$ 的交, $n\in\mathbb{Z}_+$. 对于每一个 $n$, 选取一个连续函数 $f_n\colon X\to[0,1]$, 使得当 $x\in A$ 时, $f_n(x)=0$, 当 $x\in X-U_n$ 时 $f_n(x)=1$. 定义函数 $f(x)=\sum f_n(x)/2^n$. 通过与 $\sum 1/2^n$ 比较可见, 这个序列是一致收敛的, 所以 $f$ 是连续的, 并且在 $A$ 上取值为 $0$, 在 $X-A$ 上取正值.
\end{proof}

\begin{theorem}[Nagata-Smirnov 度量化定理]\label{theorem:ch06-nagata-smirnov-metrization}
空间 $X$ 是可度量化的当且仅当 $X$ 是正则的并且有一个可数局部有限基.
\end{theorem}

\begin{proof}
第一步. 设 $X$ 是正则的并且有一个可数局部有限基 $\mathcal{B}$, 则 $X$ 是正规的, 并且 $X$ 的每一个闭子集是一个 $G_\delta$ 集. 对于某一个 $J$, 下面我们将把 $X$ 嵌入到度量空间 $(\mathbb{R}^J,\bar{\rho})$ 中, 从而说明 $X$ 是可度量化的.

设 $\mathcal{B}=\bigcup\mathcal{B}_n$, 其中每一个集族 $\mathcal{B}_n$ 是局部有限的. 对于每一个正整数 $n$ 和 $\mathcal{B}_n$ 中的任意一个基元素 $B$, 选取一个连续函数
\begin{align*}
f_{n,B}\colon X\longrightarrow[0,1/n],
\end{align*}
使得对任意 $x\in B$ 有 $f_{n,B}(x)>0$, 对于 $x\notin B$ 有 $f_{n,B}(x)=0$. 于是, 函数族 $\{f_{n,B}\}$ 分离 $X$ 中的点和闭集: 因为对于任意给定的点 $x_0$ 和 $x_0$ 的一个邻域 $U$, 都存在一个基元素 $B$, 使得 $x_0\in B\subset U$, 并且对某一个 $n$ 有 $B\in\mathcal{B}_n$, 从而 $f_{n,B}(x_0)>0$, 并且 $f_{n,B}(x_0)$ 在 $U$ 的补集上取零值.

设 $J$ 是 $\mathbb{Z}_+\times\mathcal{B}$ 的一个子集, 它由满足 $B\in\mathcal{B}_n$ 的所有偶对 $(n,B)$ 组成. 定义
\begin{align*}
F\colon X\longrightarrow[0,1]^J
\end{align*}
为
\begin{align*}
F(x)=(f_{n,B}(x))_{(n,B)\in J}.
\end{align*}

根据定理 34.2, 相对于 $[0,1]^J$ 的积拓扑而言, 映射 $F$ 是一个嵌入.

现在我们在 $[0,1]^J$ 上取由一致度量 $\bar{\rho}$ 所诱导的拓扑, 并且证明相对于这个拓扑, $F$ 也是一个嵌入. 这便是需要条件 $f_{n,B}(x)<1/n$ 的地方. 一致拓扑是比积拓扑更细(大)的, 因此, 相对于一致度量, 映射 $F$ 是单射, 并且将 $X$ 中的开集变成像空间 $Z=F(X)$ 的开集. 我们只要再证明 $F$ 连续就够了.

注意, 在 $\mathbb{R}^J$ 的子空间 $[0,1]^J$ 上, 一致度量等价于度量
\begin{align*}
\rho\bigl((x_\alpha),(y_\alpha)\bigr)=\sup\{|x_\alpha-y_\alpha|\}.
\end{align*}
为了证明连续性, 我们取定 $X$ 的一个点 $x_0$ 和一个数 $\varepsilon>0$, 并取定 $x_0$ 的一个邻域 $W$, 使得
\begin{align*}
x\in W\Longrightarrow\rho(F(x),F(x_0))<\varepsilon.
\end{align*}

暂时固定 $n$, 选取 $x_0$ 的一个邻域 $U_n$, 使得它只与族 $\mathcal{B}_n$ 中的有限多个成员相交. 这就是说, 当 $B$ 在 $\mathcal{B}_n$ 上变动时, 函数 $f_{n,B}$ 除了有限多个以外其他的在 $U_n$ 上都恒等于 $0$. 因为每一个函数 $f_{n,B}$ 是连续的, 所以我们可以选取 $x_0$ 的一个邻域 $V_n$ 包含于 $U_n$, 并且在 $V_n$ 上, 当 $B\in\mathcal{B}_n$ 时, 每一个保留下来的函数 $f_{n,B}$ 最多改变 $\varepsilon/2$.

对于每一个 $n\in\mathbb{Z}_+$ 选择 $x_0$ 的上述邻域 $V_n$, 然后选取 $N$, 使得 $1/N\leqslant\varepsilon/2$, 并且定义 $W=V_1\cap\cdots\cap V_N$. 可以断言, $W$ 就是所要求的 $x_0$ 的那个邻域. 令 $x\in W$. 如果 $n\leqslant N$, 因为函数 $f_{n,B}$ 在 $W$ 上或者恒等于 $0$, 或者最多改变 $\varepsilon/2$, 所以
\begin{align*}
|f_{n,B}(x)-f_{n,B}(x_0)|\leqslant\varepsilon/2.
\end{align*}
如果 $n>N$, 因为 $f_{n,B}$ 把 $X$ 映入 $[0,1/n]$ 中, 所以
\begin{align*}
|f_{n,B}(x)-f_{n,B}(x_0)|\leqslant 1/n<\varepsilon/2.
\end{align*}
综上, 最后得到
\begin{align*}
\rho(F(x),F(x_0))\leqslant\varepsilon/2<\varepsilon.
\end{align*}

第二步. 现在我们来证明定理的另一方面. 设 $X$ 是可度量化的. 易见 $X$ 是正则的. 以下将证明 $X$ 有可数局部有限基.

选取 $X$ 的一个度量. 给定 $m$, 用 $\mathcal{A}_m$ 表示所有半径为 $1/m$ 的开球构成的 $X$ 的开覆盖. 根据\cref{lemma:ch06-metrizable-countable-locally-finite-refinement}可见, 存在 $X$ 的一个可数局部有限的开覆盖 $\mathcal{B}_m$ 加细 $\mathcal{A}_m$. 注意 $\mathcal{B}_m$ 中的每一个元素的直径最多为 $2/m$. 令 $\mathcal{B}=\bigcup_{m\in\mathbb{Z}_+}\mathcal{B}_m$. 因为每一个 $\mathcal{B}_m$ 是可数局部有限的, 所以 $\mathcal{B}$ 也是可数局部有限的. 我们来证明 $\mathcal{B}$ 是 $X$ 的一个基.

给定 $x\in X$ 和 $\varepsilon>0$, 下面将证明存在 $B\in\mathcal{B}$, 使得 $x\in B$ 并且 $B\subset B(x,\varepsilon)$. 先选取 $m$, 使得满足 $1/m<\varepsilon/2$. 由于 $\mathcal{B}_m$ 覆盖 $X$, 可以选取 $\mathcal{B}_m$ 中的一个元素 $B$, 使得 $x\in B$. 从而由 $x\in B$ 和 $B$ 的直径最多为 $2/m<\varepsilon$ 可见 $B$ 包含于 $B(x,\varepsilon)$.
\end{proof}

\subsection*{习题}

\begin{enumerate}
\item 验证例 1 和例 2.
\item $X$ 的子集 $W$ 称为一个 $F_\sigma$ 集, 若它等于 $X$ 的可数个闭子集的并. 证明 $W$ 是 $X$ 中的一个 $F_\sigma$ 集当且仅当 $X-W$ 是 $X$ 的 $G_\delta$ 集.

[这一术语源于法语. ``F''代表``fermé'', 意为``闭集'', ``$\sigma$''代表``somme'', 意为``并''.]
\item 许多空间都有可数基, 但 $T_1$ 空间没有局部有限基, 除非它是离散的. 试证明这一命题.
\item 找一个非离散的空间, 它有可数局部有限基, 却没有可数基.
\item $X$ 的子集族 $\mathcal{A}$ 称为局部离散的(locally discrete), 若对 $X$ 中的每一个点都存在它的一个邻域最多只与 $\mathcal{A}$ 中一个元素相交. 集族 $\mathcal{B}$ 称为可数局部离散的(countably locally discrete)(或 $\sigma$ 局部离散的), 若它等于可数个局部离散的集族的并. 证明下述命题:

\textbf{定理}(Bing 可度量定理) 一个空间 $X$ 是可度量化的当且仅当它是正则的并且有一可数局部离散基.
\end{enumerate}

\end{document}